\section{Phenomenology of the Medieval Mind}

\begin{quotex}
The man who knows the future is the man who wills the future.
\end{quotex}

In \textit{The Discarded Image}, \textbf{C. S. Lewis} offers this somewhat strange suggestion:

\begin{quotex}
[Facts] become valuable only in so far as they enable us to enter more fully into the consciousness of our ancestors by realising how such a universe must have affected those who believed in it. The recipe for such realisation is not the study of books. You must go out on a starry night and walk about for half an hour trying to see the sky in terms of the old cosmology. 

\end{quotex}
If that old cosmology is false, as Lewis believes and as it seems to be from the purely quantitative aspect, what would be the point of that exercise? Lewis thinks like a modern. There is a world ``out there, right now". It is experienced by a self-enclosed mind, which then devises various theories of models to explicate that world. This is hardly the ``consciousness of our ancestors".

The Medieval mind experienced things quite differently. For him, man, situated between the worlds of spirit and matter, was thus a Microcosm of the enclosing Macrocosm. The physical world is a reflection of the spiritual world. In knowing the physical world, man simultaneously knew the spiritual world, participating in the realm of Ideas. Modern man, whose consciousness is formed by the thoughts of a Francis Bacon\footnote{\url{https://gornahoor.net/?p=3468}}, ignores the realm of ideas. Instead, he experiences phenomena and investigates the causes of their motions.

The Medieval man is ruled by the Logos, that is, Logic, for which cause and effect are meaningless. He is interested in Form, the ``what", the ``why", the ``purpose". For him, there is a correspondence between the spiritual and the material, the inner world and the outer. It is wildly wide of the mark, for example, to believe that the Medieval mind thinks that the planets affected his consciousness like some strange sort of physical force. Rather, he sees a correspondence between the starry night and his inner states, not an abstract cause. So, a better meditation would be, while looking out at the heavens, to also look within and recognize the correspondence. Below, we give some examples from The Discarded Image.

\paragraph{Wits and Senses}
Corresponding to the five outer senses of touch, smell, taste, sight, and hearing, there are the five inner wits.

\begin{table}[h]\small\centering
\begin{tabular}{lp{.7\textwidth}}\toprule
Memory &
The ability to recall past experiences\\
Estimation &
Subconscious awareness, knowing ``in the gut"\\
Imagination &
The formation of images of external things\\
Fantasy &
The ability to consider in one's mind things that have not been experienced\\
Common sense &
The ability to recognize a unity from the many impressions coming from the external senses.\\\bottomrule
\end{tabular}
\end{table}
\paragraph{Humours}
As in all Traditional and esoteric teachings, the Medievals were concerned with identifying personality types. They recognized four main types corresponding to (not caused by) the qualitative elements and bodily processes, as well as the combinations thereof. An active meditation is to recognize these qualities in those you meet during the day.

\begin{table}[h]\small\centering
\begin{tabular}{ccc}\toprule
\textbf{Element} &
\textbf{Fluid} &
\textbf{Humour}\\\midrule
Air &
Blood &
Sanguine\\
Fire &
Yellow Bile &
Choleric\\
Earth &
Black Bile &
Melancholic\\
Water &
Phlegm &
Phlegmatic\\\bottomrule
\end{tabular}
\end{table}
\paragraph{Planets}
The planets, too, correspond to various moods, few of which are even recognizable to contemporary men. Note how few of these you use in your quotidian life. Again, we note the triple division: from the spirit, we have the spiritual hierarchy\footnote{\url{http://www.gornahoor.net/?p=1941}}; from matter, the element; finally, the correspondence to an inner state.

\begin{table}[h]\small\centering
\begin{tabular}{ccc}\toprule
\textbf{Planet} &
\textbf{Metal} &
\textbf{Mood}\\\midrule
Saturn &
Lead &
Saturnine\\
Jupiter &
Tin &
Jovial\\
Mars &
Iron &
Martial\\
Sol &
Gold &
Solar\\
Venus &
Copper &
Amorous\\
Mercury &
Quicksilver &
Mercurial\\
Moon &
Silver &
Lunar\\\bottomrule
\end{tabular}
\end{table}
As Lewis rightly points out, it is insufficient to know about these moods, but rather to know them intuitively, in one's own consciousness. Of course, the motions of the planets, or the variations of our moods, appear in various constellations or patterns\footnote{\url{https://gornahoor.net/?p=2377}}. These we must learn to recognize in ourselves, lest we be controlled by them. In thus knowing ourselves, we become master of our own destinies. The man who knows the future is the man who wills the future.



\flrightit{Posted on 2011-12-18 by Cologero }

\begin{center}* * *\end{center}

\begin{footnotesize}\begin{sffamily}



\texttt{logres on 2011-12-20 at 21:11 said: }

Taste goes with common sense, estimation with smell, fantasy with hearing, memory with touch, \& imagination with sight? Need to look this one up…

\hfill

\end{sffamily}\end{footnotesize}
